\centerline{\bf \Huge K-k-kombo}
\centerline{\bf \Huge }

\begin{flushright}
        \scriptsize{Как новый год встретишь, так его и комбигеома}
\end{flushright}

\textit{Теорема Хелли для прямой.} На прямой дано конечное множество отрезков. Известно, что любые два имеют общую точку. Тогда все отрезки имеют общую точку.

\problem{1} На прямой дано $2 n+1$ отрезков. Известно, что каждый пересекает не менее $n$ других. Докажите, что существует отрезок, пересекающий все остальные.

\problem{2} Даны несколько прямоугольников с параллельными сторонами. Известно, что любые два имеют общую точку. Докажите, что все они имеют общую точку.

\problem{3} На плоскости расположено несколько равных квадратов со сторонами, параллельными осям координат так, что каждая точка покрыта не более, чем $k$ квадратами. Докажите, что квадраты можно разбить на $2 k-1$ семейств так, чтобы в каждом семействе квадраты не пересекались.

\problem{4} Пусть на прямой даны $n k+1$ отрезок. Докажите, что из этих отрезков можно выделить или $n+1$ попарно непересекающихся, или $k+1$, имеющих общую точку.

\problem{5} На плоскости дано конечное множество квадратов $A$ с параллельными сторонами. Докажите, что можно выделить подмножество $B$ множества $A$ так, чтобы квадраты из $B$ не пересекались, а также квадраты $B$, раздутые в 3 раза, покрывали $A$.

\problem{6} На плоскости дано несколько единичных квадратов с параллельными сторонами. Пусть площадь, которую они покрывают, равна $S$. Докажите, что можно выкинуть часть квадратов так, чтобы оставшиеся не пересекались и покрывали площадь, не меньшую $S / 4$.